\section{Evaluation}
\label{sec:evaluation}

This section reports the comparative results across both configurations. The Configuration A subsection establishes the optimal baseline and reads the Q-Learning and SARSA residuals against it. The Configuration B subsection presents the function-approximation results, including a bridge experiment that decomposes the wrapper cost from the observation-corruption cost.

\subsection{Configuration A results}
\label{sec:evaluation-A}

\paragraph{Convergence of VI.} Value Iteration converges in 237 iterations to a final Bellman residual below $10^{-10}$. The greedy policy stabilises far earlier: from iteration 37 onward no state's argmax action changes. The remaining 200 iterations refine the value function numerically without affecting the policy used at deployment. \Cref{fig:vi-convergence} shows both views.

\begin{figure}[!htbp]
\centering
\includegraphics[width=0.95\linewidth]{configA_vi_convergence.png}
\caption{Value Iteration convergence. Left: Bellman residual $\|V_{k+1} - V_k\|_\infty$ on a log scale. Right: number of states with a greedy-policy change per iteration.}
\label{fig:vi-convergence}
\end{figure}

\paragraph{Algorithm comparison.} \Cref{tab:configA-results} reports the four agents on the same evaluation protocol. The random baseline establishes the floor at 67.5\% survival. Value Iteration sets the model-based ceiling at 78.8\% survival, with analytical $J^* = 0.7853$ from the linear solve. Q-Learning and SARSA reach 76.1\% and 74.3\% pooled survival respectively, with per-seed standard deviations under 2 percentage points in both cases. The two algorithms are at parity to within seed noise: their analytical $J(\pi)$ values are $0.689 \pm 0.005$ for Q-Learning and $0.689 \pm 0.003$ for SARSA, indistinguishable.

\begin{table}[!htbp]
\centering
\small
\begin{tabular}{lrrrr}
\toprule
Algorithm & Analytical $J(\pi)$ & Survival \% [95\% CI] & Mean return $\pm$ std & Mean episode length \\
\midrule
Random           & $0.605$                & $67.5$ $[64.6, 70.2]$ & $0.571 \pm 0.464$ & $10.4$ \\
Q-Learning       & $0.689 \pm 0.005$      & $76.1$ $[74.5, 77.7]$ & $0.671 \pm 0.425$ & $11.0$ \\
SARSA            & $0.689 \pm 0.003$      & $74.3$ $[72.9, 75.9]$ & $0.670 \pm 0.436$ & $9.6$ \\
Value Iteration  & $0.785$                & $78.8$ $[76.4, 81.5]$ & $0.751 \pm 0.404$ & $9.8$ \\
\bottomrule
\end{tabular}
\caption{Configuration A: per-agent performance over 1{,}000 evaluation episodes. Survival shown with 95\% bootstrap confidence interval; $J(\pi)$ from VI's analytical linear solve.}
\label{tab:configA-results}
\end{table}

\paragraph{Optimality gap.} \Cref{fig:opt-gap} tracks the optimality ratio $J(\pi_t) / J(\pi_\text{VI})$ over training by computing the analytical $J$ of each snapshot policy. Both Q-Learning and SARSA rise from the random-policy ratio of 0.771 to approximately 0.85 within 25{,}000 episodes, then continue climbing slowly and plateau at $0.877 \pm 0.006$ and $0.877 \pm 0.003$ respectively after roughly 100{,}000 episodes. The residual 12.3 percentage points to the VI optimum is the joint cost of the constant step size $\alpha = 0.1$, which prevents the Q-values from fully settling onto the argmax, and the $\epsilon = 0.05$ floor, which keeps some training transitions exploratory. Closing the last gap would require an $\alpha$-schedule, not more episodes.

\begin{figure}[!htbp]
\centering
\includegraphics[width=0.85\linewidth]{configA_optimality_gap.png}
\caption{Optimality ratio $J(\pi_t)/J(\pi_\text{VI})$ over training episodes, computed analytically from VI's linear solve. Bands span $\pm 1$ standard deviation across three seeds.}
\label{fig:opt-gap}
\end{figure}

\paragraph{Exploration regime sweep.} \Cref{fig:explore-exploit} reports the exploration-design sweep over five $\epsilon$ regimes at a 50{,}000-episode budget. The ranking is unexpected: pure exploit ($\epsilon = 0$) wins at $J/J^* = 0.917$, ahead of fixed $\epsilon = 0.3$ at $0.860$, decay-to-$0.01$ at $0.834$, and pure random and decay-to-$0.05$ at $0.829$. The result inverts the textbook ordering.

\begin{figure}[!htbp]
\centering
\includegraphics[width=0.85\linewidth]{configA_explore_exploit.png}
\caption{Optimality ratio under five $\epsilon$ regimes (single seed, 50{,}000 episodes per regime).}
\label{fig:explore-exploit}
\end{figure}

The explanation runs as follows. With $Q$ initialised to zero, $\epsilon = 0$ picks action 0 = (V=none, F=none) at every state, and the TD updates learn the value of that fixed policy. The always-no-treatment policy is unusually competitive in this MDP because action 0 carries zero intensity penalty ($\text{INTENSITY}[0] = 0$) and is also the optimal action for many low-severity states. Its analytical $J = 0.720$, which is 92\% of $J^*$. Pure random updates the table but never acts on it, so its asymptote equals the random-policy ratio. The decay schedules at 50{,}000 episodes have only 10{,}000 post-decay episodes to settle and lag both. At the 200{,}000-episode budget used in the main Q-Learning run the decay schedule reaches 0.877; even that figure remains below the 50{,}000-episode pure-exploit ratio. The exploration / exploitation trade-off in this MDP is therefore not the textbook one: a degenerate-but-cheap policy is competitive enough that an $\epsilon$-greedy schedule with a constant step size struggles to outrun it.

\paragraph{Return distribution.} \Cref{fig:configA-violin} shows the per-agent return distributions on the 1{,}000-episode evaluation set. All four are bimodal: a death lobe just below zero (the $\lambda$-intensity penalty pulls it slightly negative) and a survival lobe near $+1$. The mean for each agent sits between the two lobes; the upward shift across Random $\to$ SARSA / Q-Learning $\to$ VI reflects mass moving from the death lobe to the survival lobe.

\begin{figure}[!htbp]
\centering
\includegraphics[width=0.85\linewidth]{configA_return_distribution.png}
\caption{Per-episode return distributions for the four Configuration A agents (1{,}000 evaluation episodes each).}
\label{fig:configA-violin}
\end{figure}

\paragraph{Policy interpretation.} \Cref{fig:configA-policies} shows the visit-weighted action heatmaps for each policy. Value Iteration is sharply concentrated: 95\% of its visits sit in the F0 (no-fluid) column, with a single dominant cell at (V0, F0) carrying 38.4\%. The optimal model-based policy is "titrate vasopressor while keeping IV fluid off," consistent with the $\lambda$-penalty making fluid expensive when it is not necessary. Q-Learning and SARSA do not recover this sparsity. Both spread mass across the whole $5 \times 5$ grid: Q-Learning's largest cell is (V4, F0) at 12.5\%, SARSA's is (V1, F0) at 13.1\%, with nothing approaching VI's 38\%. The diffuse pattern reflects the constant $\alpha = 0.1$ and $\epsilon$ floor: near-tied Q-values across actions are not resolved into a sharp argmax. The same effect explains why both algorithms also show essentially no SOFA-severity dependence in their action distributions, whereas VI mildly escalates vasopressor with severity.

\begin{figure}[!htbp]
\centering
\includegraphics[width=\linewidth]{configA_policy_actions.png}
\caption{Visit-weighted greedy-action heatmaps. Rows: vasopressor level (V0--V4); columns: IV fluid level (F0--F4).}
\label{fig:configA-policies}
\end{figure}

\paragraph{Saveability across states.} The optimal value function $V^*(s)$ falls monotonically with SOFA severity, from approximately 0.90 for SOFA $< 5$ down to 0.20 for the sickest patients. The histogram of $V^*$ across the 714 non-terminal states has a heavy ridge at $V^* \approx 0.85$ ($\sim 700$ states) and a long left tail with roughly 30 states below $V^* < 0.5$. The analytical $J^* = d_0 \cdot V^* = 0.785$ sits below the unweighted state-average $V^* = 0.826$ because the SOFA-biased initial distribution puts mass on harder states. The 78.8\% survival ceiling is therefore not a learning artefact: it is the model-based optimum given the cohort, and the 12 percentage-point gap between unweighted mean $V^*$ and $J^*$ quantifies the cohort difficulty introduced by SOFA bias.

\begin{figure}[!htbp]
\centering
\includegraphics[width=0.95\linewidth]{configA_value_landscape.png}
\caption{Left: $V^*(s)$ vs SOFA score, marker size $\propto d_0(s)$. Right: histogram of $V^*$ across the 714 non-terminal states.}
\label{fig:value-landscape}
\end{figure}

\subsection{Configuration B results}
\label{sec:evaluation-B}

\paragraph{Random baseline on the clinical environment.} The clinical environment lowers the random-policy survival floor from 67.5\% (Configuration A) to 65.0\% (\Cref{fig:configB-floor}). The 2.5 percentage-point drop is entirely the acute-event wrapper: when an acute event fires (7.9\% of episodes), the patient dies regardless of action, so survival is 0\% in that subset. Removing those episodes lifts the random survival to 70.6\%. The observation-noise and missing-lab wrappers barely move the random baseline (71.5\% and 69.3\% in their respective subsets, within sampling noise of overall), because a random policy never looks at the observation. For learning agents the noise and missing wrappers do matter, but the random floor isolates the acute-event tax cleanly.

\begin{figure}[!htbp]
\centering
\includegraphics[width=0.85\linewidth]{configB_env_exploration.png}
\caption{Random-policy survival on the clinical environment, decomposed by failure-mode subset (1{,}000 episodes).}
\label{fig:configB-floor}
\end{figure}

\paragraph{Both Linear-Q variants collapse to a single action.} Naive Linear-Q's training curve climbs from the random baseline (0.559) to 0.685 within 40{,}000 environment steps and stays perfectly flat thereafter. All three seeds collapse onto an identical asymptote. RFF Linear-Q oscillates between 0.59 and 0.68 across the budget; the across-seed band is wide and unstable. Both deploy to nearly identical policies in the action heatmap (\Cref{fig:configB-policies}): naive Linear-Q places 100\% of visits at (V=none, F=none); RFF Linear-Q places 96\%, with the remaining 4\% spread across unrelated cells that do not respond to severity.

This is not a tie-breaking artefact. Weights are initialised from $\mathcal{N}(0, 0.01^2)$ to break argmax ties, and three seeds with different initial weights converge to the same attractor for both representations. The RFF result refines the initial hypothesis. Linear-Q's failure is not representational in the narrow sense of "the linear functional form cannot express the optimum"; even a 5x richer linear-in-features class (256-dim RBF approximation) still collapses. The basin of always-action-0 is a structural property of the MDP's loss landscape under TD. Section \ref{sec:evaluation-A} provides the analytical justification: the always-action-0 policy reaches 92\% of $J^*$ on the clean MDP, so the gradient signal pointing away from it is weak regardless of the feature space.

\paragraph{DQN learns a non-degenerate policy with reduced oscillation after tuning.} The tuned DQN curve oscillates within roughly 0.55 to 0.65 across the 300{,}000-step budget, but the dynamics are visibly smoother than the original 150{,}000-step run at default hyperparameters. The mean snapshot return shows a clear upward trend after the first 100{,}000 steps and ends at $\sim 0.62$ across three seeds. Soft target updates and the five-times-smaller learning rate together remove a major source of training instability. The deployed best-snapshot model reaches 68.1\% survival and mean return 0.640, compared to 67.4\% and 0.608 at the original 150{,}000-step setup.

\Cref{tab:configB-results} reports the Configuration B comparison. Mean returns rank Linear-Q-raw (0.666) above Linear-Q-RFF (0.665) above DQN (0.640) above random (0.559); survival ranks DQN (68.1\%) above Linear-Q-RFF (66.7\%) above Linear-Q-raw (66.6\%) above random (65.0\%). The two Linear-Q variants are statistically indistinguishable on both metrics; the deep approximator is the only one that meaningfully exceeds the linear floor.

\begin{table}[!htbp]
\centering
\small
\begin{tabular}{lrrrr}
\toprule
Algorithm & Mean return $\pm$ std & Survival \% [95\% CI] & Mean episode length & Convergence to 90\% of best \\
\midrule
Random            & $0.559 \pm 0.477$ & $65.0$ $[61.9, 68.2]$ & $9.2$ & n/a \\
Linear-Q (raw)    & $0.666 \pm 0.472$ & $66.6$ $[63.7, 69.7]$ & $9.3$ & $40{,}000$ steps \\
Linear-Q (RFF)    & $0.665 \pm 0.471$ & $66.7$ $[63.6, 69.5]$ & $8.8$ & $150{,}000$ steps \\
DQN               & $0.640 \pm 0.467$ & $68.1$ $[65.2, 70.8]$ & $9.3$ & not reached \\
\bottomrule
\end{tabular}
\caption{Configuration B: per-agent performance on the clinical environment (1{,}000 evaluation episodes, fixed seed 42).}
\label{tab:configB-results}
\end{table}

\begin{figure}[!htbp]
\centering
\includegraphics[width=\linewidth]{configB_learning_curves.png}
\caption{Configuration B learning curves: mean snapshot return across three seeds with $\pm 1$ standard deviation band.}
\label{fig:configB-curves}
\end{figure}

The mismatch is explained by the $\lambda$-intensity penalty. Linear-Q's always-action-0 policy never incurs any intensity cost, so when the patient dies Linear-Q's return is exactly 0 rather than below 0. DQN treats the patient and accumulates intensity cost, pulling its lower-lobe returns below zero even though more patients survive. Mean return rewards Linear-Q's cost avoidance; survival rewards DQN's actual behaviour. Survival is the correct metric for this clinical task; mean return is contaminated by the $\lambda$ design.

\begin{figure}[!htbp]
\centering
\includegraphics[width=\linewidth]{configB_policy_actions.png}
\caption{Configuration B greedy-action frequency heatmaps. Rows: vasopressor level; columns: IV fluid level.}
\label{fig:configB-policies}
\end{figure}

\paragraph{Clinical interpretability.} Tuned DQN is the only Configuration B agent that responds to severity. Mean vasopressor level rises from approximately $0.75$ at start-SOFA $\approx 2$ to $1.44$ at start-SOFA $\approx 18$, a monotonic $\sim 0.7$-unit escalation on the 0--4 scale (\Cref{fig:configB-sofa}). Mean fluid level rises from $\sim 0.5$ at low SOFA to a peak of $\sim 1.0$ at SOFA $\approx 10$, then plateaus around $0.85$ at higher SOFA --- so vasopressor moves monotonically with severity, fluid moves less and not monotonically. Both Linear-Q variants (raw and RFF) are essentially flat at 0 across the SOFA axis for both treatments, consistent with their single-action policies. DQN's escalation pattern is qualitatively the one Value Iteration showed in Configuration A: vasopressor moves with severity, fluid moves less. The magnitude is moderate, which is consistent with DQN's +3.1 percentage-point survival lift over the random floor.

\begin{figure}[!htbp]
\centering
\includegraphics[width=\linewidth]{configB_action_by_sofa.png}
\caption{Mean treatment level by start-SOFA bin: vasopressor (left) and IV fluid (right).}
\label{fig:configB-sofa}
\end{figure}

\subsection{Bridge experiment: tabular policy on the clinical environment}
\label{sec:evaluation-bridge}

A direct survival-number comparison between Configuration A and Configuration B conflates two effects: the change in observation regime, and the addition of the clinical wrappers. To separate them, this experiment deploys Configuration A's optimal Value Iteration policy $\pi_\text{VI}$ on the Configuration B clinical environment. The deployment reads the underlying discrete state index at each step (peeking past the observation wrappers) and selects $\pi_\text{VI}(s)$. The agent therefore has perfect state information but still suffers irreducible acute-event mortality.

The oracle achieves 78.5\% survival on the clinical environment, only 0.3 percentage points below VI's 78.8\% on the clean environment. The wrappers cost the optimal tabular policy almost nothing because the noise and missing-lab wrappers do not change the action of a policy that is indexed on the underlying state, and the acute-event wrapper subtracts only its 0.3 percentage-point irreducible tax.

\Cref{tab:comparison-AB} and \Cref{fig:comparison} summarise the full comparison.

\begin{table}[!htbp]
\centering
\small
\begin{tabular}{llrr}
\toprule
Setting & Method class & Mean return & Survival \% \\
\midrule
VI (Config A, clean discrete state)       & tabular DP                    & $0.751$ & $78.8$ \\
VI policy on clinical env (oracle)        & tabular policy, deployed       & $0.749$ & $78.5$ \\
Q-Learning policy on clinical env (oracle) & tabular policy, deployed       & $0.645$ & $72.3$ \\
SARSA policy on clinical env (oracle)      & tabular policy, deployed       & $0.640$ & $70.4$ \\
Linear-Q (Config B)                       & linear function approximation  & $0.666$ & $66.6$ \\
DQN (Config B)                            & deep function approximation    & $0.640$ & $68.1$ \\
\bottomrule
\end{tabular}
\caption{Configuration A vs B comparison. The four oracle rows allow a three-way decomposition: wrapper cost ($\sim 0.3$ pp), algorithm-quality gap ($\sim 6$ pp from VI to Q-Learning tabular policy on the same clinical env), and observation-regime cost ($\sim 5$ pp from Q-Learning tabular to DQN function approximation).}
\label{tab:comparison-AB}
\end{table}

\begin{figure}[!htbp]
\centering
\includegraphics[width=\linewidth]{comparison_AB.png}
\caption{End-to-end comparison. Left: survival rate. Right: mean return.}
\label{fig:comparison}
\end{figure}

The expanded bridge experiment decomposes the gap from VI on clean to tuned DQN on clinical (10.7 pp) into three measurable parts. The wrapper cost is 0.3 pp, all of it the irreducible acute-event tax: VI's sparse "no-fluid" policy is robust to observation corruption because its action choices do not depend on the noisy observations the wrappers modify. The algorithm-quality gap is 6.2 pp: Q-Learning's tabular policy deployed on the clinical env reaches only 72.3\% survival, 6.2 pp below the VI oracle on the same env. The observation-regime cost is 4.2 pp: tuned DQN at 68.1\% sits below Q-Learning tabular at 72.3\% by exactly that gap, despite running on the same MDP. Q-Learning and SARSA tabular policies lose 4.5 and 3.9 pp respectively going from clean to clinical, an order of magnitude more than VI loses. Their diffuse action distributions interact differently with episodic termination than VI's sparse policy does.

The decomposition tightens the original argument. The clinical wrappers are not what the Configuration B agents are struggling against. The dominant Configuration B cost is split between the algorithm-quality gap (which would exist on any policy below VI's optimum) and the observation regime (which is unavoidable when only raw features are available). Algorithm-class choice should follow the observation model, but even with perfect tabular policies, the wrapper-induced robustness gap (4--5 pp for model-free TD vs 0.3 pp for VI) is itself a non-trivial cost worth measuring.
